\documentclass[11pt,a4paper]{article}
\usepackage[utf8]{inputenc}
\usepackage[margin=2.5cm]{geometry}
\usepackage{amsmath, amssymb}
\usepackage{booktabs}
\usepackage{hyperref}

\hypersetup{
    colorlinks=true,
    linkcolor=blue,
    urlcolor=cyan,
    pdftitle={Appunti Tex-Notes}
}

\title{\textbf{Tex-Notes}: Appunti Rapidi in \LaTeX}
\author{Autore del Documento}
\date{\today}

\begin{document}

\maketitle

\section{Introduzione}
Benvenuto in \textbf{Tex-Notes}! Questo editor minimale e completamente offline ti permette di scrivere appunti in \LaTeX{} con anteprima live ultra-veloce e consultazione integrata dei comandi.

\subsection{Caratteristiche Principali}
\begin{itemize}
    \item \textbf{Compilazione Offline:} Tectonic compila localmente senza richiedere connessione internet.
    \item \textbf{Anteprima Live:} Aggiornamento automatico dopo 1-2 secondi di pausa nella digitazione.
    \item \textbf{Vocabolario dei Comandi:} Clicca su un comando dal pannello laterale per inserirlo al cursore.
\end{itemize}

\section{Matematica ed Equazioni}
Possiamo inserire formule in linea come $e^{i\pi} + 1 = 0$ oppure formule centrate:

\begin{equation}
    f(x) = \int_{0}^{\infty} \frac{\sin(t)}{t} \, dt = \frac{\pi}{2}
\end{equation}

\subsection{Matrici ed Algebra Lineare}
Ecco una matrice tra parentesi tonde e una tra parentesi quadre:
\begin{equation*}
    A = \begin{pmatrix}
        \cos\theta & -\sin\theta \\
        \sin\theta &  \cos\theta
    \end{pmatrix}, \quad
    B = \begin{bmatrix}
        1 & 2 & 3 \\
        0 & 1 & 4 \\
        5 & 6 & 0
    \end{bmatrix}
\end{equation*}

\section{Tabelle Eleganti}
Grazie al pacchetto \texttt{booktabs}, le tabelle risultano pulite e professionali:

\begin{table}[htbp]
\centering
\begin{tabular}{llr}
\toprule
\textbf{Comando} & \textbf{Descrizione} & \textbf{Esempio} \\
\midrule
\texttt{\textbackslash frac\{a\}\{b\}} & Frazione & $\frac{1}{2}$ \\
\texttt{\textbackslash sum\_\{i=1\}\^{}n} & Sommatoria & $\sum_{i=1}^n i$ \\
\texttt{\textbackslash textbf\{testo\}} & Grassetto & \textbf{testo} \\
\bottomrule
\end{tabular}
\caption{Esempio di comandi rapidi.}
\label{tab:comandi}
\end{table}

\end{document}
